\originalchapter{8.05 2011测试：变分与算符}
原件为Fall 2011 Test, 2011年10月26日12:30--14:00, 五题100分, 闭卷禁计算器; 由Fall 2013 OCW提供. 以下为\textbf{AI独立解答}, 判断题加简短理由帮助复习, 原卷仅要求T/F.
\originalsection{第1题：十个判断（20分）}
\begin{exercise}
前三项势为$x<0$的$2V_0$、$0\leq x\leq a$的0、$x>a$的$V_0$, $V_0>0$. 第(4)、(5)项自旋为$1/2$. 判断: (1)$V_0<E<2V_0$每能量有两个本征态; (2)$0<E<V_0$每能量有一个正规态; (3)$E<0$没有能量本征态; (4)$z$向上自旋与$x$向上正交; (5)$S_x^2,S_y^2$对易; (6)两幺正积幺正; (7)投影可有本征值$-1$; (8)能同时对角化的两算符对易; (9)Hermitian算符指数幺正; (10)两Hermitian算符对易子Hermitian.
\end{exercise}
\begin{solution}
依次为\textbf{F,F,T,F,T,T,F,T,F,F}. (1)只有右侧一个传播通道, 左侧衰减条件与匹配留一个散射解. (2)束缚能离散, 不能每个区间能量都有态. (3)$H$非负. (4)重叠幅$1/\sqrt2$. (5)自旋$1/2$的每个平方都是$\hbar^2I/4$. (6)$(UV)^\dagger UV=I$. (7)$P^2=P$使本征值仅0或1. (8)共同基中对角矩阵积可换序. (9)实本征值指数为正数, 一般不是模1; $\exp(\ii A)$才幺正. (10)$[A,B]^\dagger=-[A,B]$, 是反Hermitian（零是特例）.
\end{solution}
\originalsection{第2题：迹与共同基（20分）}
\begin{exercise}
(a)证明$\Tr M=\sum_i\bra{a_i}M\ket{a_i}$与正交归一基无关. (b)给$A_1=\left(\begin{smallmatrix}1&0&1\\0&0&0\\1&0&1\end{smallmatrix}\right)$,$A_2=\left(\begin{smallmatrix}2&1&1\\1&0&-1\\1&-1&2\end{smallmatrix}\right)$. 在$u_1=(1,0,1)/\sqrt2,u_2=(-1,0,1)/\sqrt2,u_3=(0,1,0)$基中补全$A_2$并求其谱.
\end{exercise}
\begin{solution}
有限维时所有和有限; 无穷维须限制迹类算符及合法求和, 不能直接换序任意发散级数. 插入另一完备基得$\sum_i\sum_{jk}\braket{b_i}{a_j}\bra{a_j}M\ket{a_k}\braket{a_k}{b_i}=\sum_jM_{jj}$, 因中间求和为$\delta_{kj}$. 有限矩阵也即$\Tr(U^\dagger MU)=\Tr M$. 逐列算$A_2u_1=3u_1$,$A_2u_2=u_2-\sqrt2u_3$,$A_2u_3=-\sqrt2u_2$, 所以矩阵为$\left(\begin{smallmatrix}3&0&0\\0&1&-\sqrt2\\0&-\sqrt2&0\end{smallmatrix}\right)$, 剩下二阶特征式$\lambda^2-\lambda-2=0$, 谱$3,2,-1$. $A_1$零特征值空间内还要对角化$A_2$, 不能认为给出的每个$A_1$本征矢已经是共同矢.
\end{solution}
\originalsection{第3题：线性绝对值势的上界（20分）}
\begin{exercise}
$V=V_0|x|$, $V_0$量纲为能量/长度, 用$\psi=A\e^{-bx^2/2}$,$b>0$优化基态上界, 化为纯数乘$\hbar,m,V_0$构成的尺度.
\end{exercise}
\begin{solution}
归一后$A=(b/\pi)^{1/4}$,$\langle T\rangle=\hbar^2b/(4m)$,$\langle|x|\rangle=1/\sqrt{\pi b}$. 因而$F(b)=\hbar^2b/(4m)+V_0/\sqrt{\pi b}$. $F'=0$给$b_*=[2mV_0/(\sqrt\pi\hbar^2)]^{2/3}$, 两项在该点势能为动能2倍. 上界
\[
E_g\leq F(b_*)=\frac3{2^{4/3}\pi^{1/3}}\left(\frac{\hbar^2V_0^2}{m}\right)^{1/3}.
\]
该尺度由动能约$\hbar^2/(mL^2)$与势能$V_0L$平衡得到; 变分固定纯数使“量级”变成真正上界.
\end{solution}
\originalsection{第4题：八次势的下界（20分）}
\begin{exercise}
$V=\gamma x^8$, $\gamma>0$. (a)将基态$\langle x^8\rangle$与$\Delta x$联系. (b)用不确定性给基态能量下界并化简尺度.
\end{exercise}
\begin{solution}
偶势一维基态唯一且偶, $\langle x\rangle=0$. 连续两次Cauchy--Schwarz给$\langle x^8\rangle\geq\langle x^4\rangle^2$, $\langle x^4\rangle\geq\langle x^2\rangle^2$, 合起来给$\langle x^8\rangle\geq\langle x^2\rangle^4=(\Delta x)^8$. $\langle p\rangle=0$, 不确定性得$\langle T\rangle\geq\hbar^2/[8m(\Delta x)^2]$. 令$s=\Delta x,A=\hbar^2/(8m)$, $E_g\geq A/s^2+\gamma s^8$. 在正$s$求最小:$s^{10}=A/(4\gamma)$, 动能界项为势能界项4倍, 得
\[
E_g\geq5\,4^{-4/5}\left(\frac{\hbar^2}{8m}\right)^{4/5}\gamma^{1/5}.
\]
这里是在所有允许宽度中最小化一个\textbf{下界}, 与试探态能量的上界不同. 两个不等式不能同时饱和于真实八次势基态, 所以无需期望等号.
\end{solution}
\originalsection{第5题：绕$y$轴进动（20分）}
\begin{exercise}
自旋$1/2$的$H=-\gamma BS_y$, 初方向$\theta_0,\phi_0=0$. (a)写初态. (b)求演化. (c)给方向角函数.
\end{exercise}
\begin{solution}
初态$\cos(\theta_0/2)\ket++\sin(\theta_0/2)\ket-$. 令$b=\gamma Bt$, $U=\cos(b/2)I+\ii\sin(b/2)\sigma_y$, 其矩阵为$\left(\begin{smallmatrix}\cos(b/2)&\sin(b/2)\\-\sin(b/2)&\cos(b/2)\end{smallmatrix}\right)$. 演化为$\cos[(\theta_0-b)/2]\ket++\sin[(\theta_0-b)/2]\ket-$. 方向$\boldsymbol n=(\sin(\theta_0-b),0,\cos(\theta_0-b))$. 标准角$\theta=\arccos\cos(\theta_0-b)$, 当$x$分量非负取$\phi=0$, 负取$\phi=\pi$; 极点$\phi$无意义. 直接让$\theta_0-b$跨过负值而仍叫标准极角会混淆, 向量表达没有这个问题.
\end{solution}
